\SourceLeaf{362}binette im J. 1833 gezogene Kette zu durchlaufen, an welcher allein der in der Luft befindliche Draht eine Länge von fast 7000 Fuss hat. Der Hauptzweck dieser Anlage ist zwar, physikalische Untersuchungen über die Gesetze der galvanischen Ströme im grossen Maassstabe anzustellen; aber gleich von Anfang an war die Gelegenheit auch vielfältig zu Versuchen einer electromagnetischen Telegraphie benutzt, die auch mit ganzen Wörtern und kleinen Phrasen auf das befriedigendste gelangen, wie schon in diesen Blättern, bei Gelegenheit der ersten Nachricht von der Einrichtung des hiesigen magnetischen Observatoriums erwähnt ist (Gött. gel. Anz. 1834, Aug. 9). An die Stelle des dabei zuerst angewandten Verfahrens wurde 1835 ein anderes gesetzt, auf welches der Hofr. \textsc{Gauss} durch die Erwägung der Inductionsgesetze geführt war, und welches dem zuerst gebrauchten bei weitem vorzuziehen ist. Gerade bei dieser Art des Telegraphirens hat nun auch der neue in Rede stehende Apparat einen bedeutenden Vorzug vor dem gewöhnlichen Magnetometer, und von diesem Umstande nahm der Hofr. \textsc{Gauss} Veranlassung, dieses Verfahren, welches bisher noch nicht öffentlich erwähnt war, in der Vorlesung nach seinen Hauptzügen zu beschreiben, und was dasselbe leistet, anzugeben, was wir jedoch hier, des beschränkten Raumes wegen, mit Stillschweigen übergehen müssen. Aus demselben Grunde erwähnen wir hier auch nur kurz einer andern neuen Vorrichtung, die zum Zwecke hat, jede unzeitige, bei einem bestimmten Geschäft störende, Schwingungsbewegung einer Magnetnadel in kurzer Zeit von selbst zur Ruhe zu bringen, und die deshalb ein Dämpfer genannt ist. Diese Vorrichtung kann eben so gut bei dem neuen Apparate wie bei dem Magnetometer gebraucht werden, und ist so wohl bei der erwähnten Art des Telegraphirens, wie bei vielen anderen magnetischen Messungsgeschäften von wesentlichem Nutzen.

\expandafter\def\csname HistoricalNoteText363\endcsname{Dieser Aufsatz enthält den wesentlichen Inhalt der in der öffentlichen Sitzung der Königlichen Societät der Wissenschaften am 19. September 1837 von mir gehaltenen Vorlesung.}
\clearpage
\fancyhead[L]{\small ÜBER EIN NEUES INSTRUMENT.}
\SourceLeaf{363}
\pdfbookmark[0]{Über ein neues Instrument. 1837}{instrument-1837}
\begin{center}
{\large ÜBER EIN NEUES, ZUNÄCHST ZUR UNMITTELBAREN BEOBACHTUNG DER VERÄNDERUNGEN IN DER INTENSITÄT DES HORIZONTALEN THEILS DES ERDMAGNETISMUS BESTIMMTES INSTRUMENT\HistoricalNote{363}.}
\par\medskip
Resultate aus den Beobachtungen des magnetischen Vereins. 1837. I.
\end{center}

Zur vollständigen Bestimmung des Erdmagnetismus an einem gegebenen Orte ist bekanntlich ein System von \textit{drei} Elementen erforderlich, und gewöhnlich wählt man dazu die Abweichung, die Neigung und die Stärke; indessen obgleich diese Wahl die für den Begriff einfachste ist, so ist es doch nicht nur verstattet, sondern es kann auch in manchen Beziehungen empfehlenswerther sein, eine andere Combination zum Grunde zu legen. Namentlich ist es sowohl in praktischer als in theoretischer Hinsicht weit vortheilhafter, den horizontalen Theil der erdmagnetischen Kraft für sich zu betrachten, und in zwei Elementen darzustellen, der Richtung (Declination) und der Stärke. Verbindet man dann damit als drittes Element entweder die Stärke der verticalen Kraft, oder die Neigung der Ganzen, so ergibt sich daraus die Stärke der ganzen Kraft, wenn man sie verlangt, von selbst.

Was nun die beiden Elemente des horizontalen Erdmagnetismus, von welchem allein hier die Rede sein wird, betrifft, so sind für die Declination durch das seit fünf Jahren eingeführte Magnetometer alle vorkommenden Aufgaben voll\SourceLeaf{364}kommen gelöst. Nicht allein zur Bestimmung ihres absoluten Werthes, sondern auch zur Verfolgung ihrer regelmässigen und zufälligen Aenderungen, von Jahr zu Jahr, von Monat zu Monat, von Stunde zu Stunde, ja selbst von einer Minute zur andern, dient dasselbe mit einer Sicherheit, Bequemlichkeit und Schärfe, die nichts zu wünschen übrig lassen.

Dasselbe Instrument dient nun zwar zugleich zur Bestimmung der Stärke des horizontalen Erdmagnetismus in absolutem Maass; ja, gerade diese Aufgabe hat, wie bekannt ist, zur Einrichtung des Magnetometers den ersten Anlass gegeben: gleichwohl löst dasselbe die Aufgabe noch keinesweges vollständig in \textit{allen} Beziehungen.

Um das, was dabei noch zu wünschen bleibt, gehörig ins Licht zu setzen, muss ich zuvörderst in Erinnerung bringen, dass die Anwendung des Magnetometers zur Bestimmung der magnetischen Intensität auf einer Verbindung \textit{mehrerer} Operationen beruht, deren Eine in der Beobachtung der Schwingungsdauer einer Nadel besteht. Diese erfordert aber ihrer Natur nach eine nicht unbeträchtliche Zeit, da die Anzahl der Schwingungen, aus denen man auf die Dauer Einer zurückschliessen muss, nicht zu klein sein darf. Ist nun während der Dauer einer solchen Operation die Intensität des Magnetismus constant, so entspricht allerdings die gefundene Schwingungsdauer \textit{diesem} Werthe der Intensität; hingegen wird jene nur \textit{dem Mittelwerthe} der Intensität während jenes Zeitraumes entsprechen, wenn dieselbe inzwischen veränderlich gewesen ist. Es bleibt uns aber auf diese Weise gänzlich verborgen, ob und was für Veränderungen in der magnetischen Intensität \textit{während} dieser Zeit vorgegangen sind. Man sieht also, dass dieses Instrument nur Durchschnittswerthe während gewisser Zeiträume geben kann, nicht aber den treuen vollständigen Hergang innerhalb derselben; wollte man, um sich diesem mehr zu nähern, die Zeiträume kürzer wählen, oder die Resultate immer nur auf eine kleine Anzahl von Schwingungen gründen, so würden jene dadurch zu sehr an Schärfe und Sicherheit verlieren, und man würde Gefahr laufen, für Anomalien in der Intensität zu halten, was nur Fehler der Beobachtungen wäre.

Je interessanter nun aber gerade die in kurzen Zeitfristen wechselnden Störungen der erdmagnetischen Kraft schon in ihrer einseitigen Erscheinung bei der Declination durch die Erfahrungen der letzten Jahre hervorgetreten sind, desto lebhafter musste man wünschen, ein Mittel zu besitzen, wodurch auch die nicht
\SourceLeaf{365}zu bezweifelnden Wirkungen solcher Störungen auf die Intensität mit derselben Leichtigkeit, Sicherheit und Schärfe verfolgt und gemessen werden könnten.,

Die Untauglichkeit der bisherigen Beobachtungsmittel zu \textit{diesem} Zwecke beruht nach dem, was ich eben entwickelt habe, darauf, dass sie auf Beobachtungen von Schwingungszeiten basirt sind, die ihrer Natur nach jedesmal eine zu lange Zeit erfordern. Die Schwingungsdauer einer Nadel dient hier aber selbst nur dazu, mittelbarerweise das Drehungsmoment zu bestimmen, welches die erdmagnetische Kraft der Nadel ertheilt, wenn sie sich nicht im magnetischen Meridian befindet. Kann man also dieses Drehungsmoment auf directem Wege, ohne Schwingungsbeobachtungen, scharf bestimmen, und seine Veränderungen sicher, scharf und schnell messen, so wird unsere Aufgabe in der Hauptsache gelöst sein. Das von mir dazu angewandte Mittel beruht auf folgender Grundlage.

Die Bedingungen des Gleichgewichts eines an zwei Fäden aufgehängten Körpers von beliebiger Gestalt, dessen Theile einstweilen bloss der Schwerkraft unterworfen und in festem Zusammenhange vorausgesetzt werden, lassen sich kurz so zusammenfassen, dass die Vertikale durch den Schwerpunkt des Körpers und die durch die Fäden dargestellten geraden Linien sich in Einer Ebene befinden, und zugleich entweder unter sich parallel sein, oder sich in Einem Punkt schneiden müssen. Allemal sind also bei der Gleichgewichtsstellung die beiden Fäden und der Schwerpunkt in Einer Vertikalebene. Um die Vorstellungen zu fixiren, mag man annehmen, dass die beiden Fäden gleich lang, ihre obern Anknüpfungspunkte in gleicher Höhe sind und von einander eben so weit abstehen, wie die beiden untern, endlich dass die letztern mit dem Schwerpunkte ein gleichschenkliges Dreieck bilden. Unter diesen Voraussetzungen werden also im Gleichgewichtszustande die beiden Fäden vertikal hängen, und eine dritte Vertikallinie, mitten zwischen diesen Fäden gedacht, wird den Schwerpunkt des Körpers treffen. Bringt man den Körper aus dieser Lage vermittelst einer Drehung um letztere Linie, so werden die beiden Fäden nicht mehr vertikal und auch nicht mehr in Einer Ebene sein, und zugleich wird der Körper etwas gehoben. Es entsteht demnach ein Bestreben, zu der vorigen Lage zurückzukehren, mit einem Drehungsmomente, welches mit hinlänglicher Genauigkeit dem Sinus der Ablenkung von der Ruhestellung proportional gesetzt werden kann, also am grössten ist, wenn die Ablenkung 90 Grad beträgt: dieses grösste Drehungsmoment wird im\SourceLeaf{366}mer stillschweigends verstanden, wenn man von Drehungsmoment schlechthin spricht. Man kann dasselbe auch als das Maass einer Kraft ansehen, mit welcher der Körper vermöge der Aufhängungsart in seiner Gleichgewichtsstellung zurückgehalten wird, und die ich der Kürze wegen die aus der Aufhängungsart entspringende Directionskraft nennen will. Ihre Grösse hängt übrigens ab 1) von der Länge der Aufhängungsfäden, 2) deren Abstande, 3) dem Gewicht des Körpers, und zwar so, dass sie der Länge der Fäden verkehrt, dem Quadrate ihres Abstandes direct, und dem Gewicht des Körpers gleichfalls direct proportional ist. Wenn die obigen Voraussetzungen nicht genau zutreffen, so ist der Ausdruck für die Directionskraft complicirter, so wie auch die Reaction der Fäden gegen eine Torsion noch eine kleine Modification nöthig macht. Es fehlt jedoch nicht an Mitteln, die Grösse der Directionskraft in grösster Schärfe durch Versuche zu bestimmen. Ueberlässt man den Körper, nach einer kleineren oder grösseren Ablenkung von der Gleichgewichtsstellung, sich selbst, so wird er mit der grössten Regelmässigkeit Schwingungen machen, deren Mitte mit dieser Stellung zusammenfällt, und deren Dauer von der Grösse der Directionskraft und dem Trägheitsmoment des Körpers abhängt.

Gehen wir jetzt zu der Voraussetzung über, dass ein horizontaler Magnetstab einen Bestandtheil des aufgehängten Körpers ausmache, so tritt eine zweite Directionskraft mit ins Spiel, und die Erscheinungen hängen von der Zusammensetzung der beiden Directionskräfte, nach den bekannten Regeln der Statik ab. Es sind in dieser Beziehung drei Fälle zu unterscheiden, indem die beiden Stellungen des Körpers, in welchen er vermöge jeder der beiden Kräfte für sich allein im Gleichgewichtszustande sein würde, entweder zusammenfallen, oder entgegengesetzt sein, oder einen Winkel mit einander machen können. Man sieht leicht, dass der Unterschied dieser drei Fälle auf dem Verhältniss der beiden Winkel beruht, welche einerseits die gerade Linie durch die beiden untern Anknüpfungspunkte der Fäden mit dem Magnetstabe, und andererseits die gerade Linie durch die beiden obern Befestigungspunkte mit dem magnetischen Meridian macht. Denkt man sich den Körper in derjenigen Gleichgewichtslage, die durch die Aufhängungsart allein bedingt wird, so wird für den ersten unsrer drei Fälle der Magnetstab im magnetischen Meridian sein müssen, und zwar in seiner natürlichen Lage (Nordpol auf der Nordseite); für den zweiten Fall muss er in verkehrter Lage im Meridian sein, und für den dritten muss er mit dem magne\SourceLeaf{367}tischen Meridian einen Winkel machen. Der Kürze wegen will ich diese drei möglichen Lagen des Magnetstabs in dem Apparate die natürliche, die verkehrte und die transversale nennen.

Bei der natürlichen Lage wird durch die Einwirkung des Erdmagnetismus auf den Magnetstab die der Aufhängungsart entsprechende Gleichgewichtsstellung des Apparats nicht abgeändert, aber dieser mit einer verstärkten Kraft darin zurückgehalten, welche die Summe der beiden Directionskräfte ist.

Im zweiten Falle, der verkehrten Lage, hört zwar das Gleichgewicht in jener Stellung auch nicht auf, allein es ist nur dann stabil, wenn die magnetische Directionskraft kleiner ist als die Directionskraft vermöge der Aufhängungsweise, und der Apparat wird dann in dieser Stellung nur mit einer Kraft zurückgehalten, die die Differenz jener beiden Directionskräfte ist. Wäre hingegen umgekehrt die magnetische Directionskraft die grössere, so würde jenes Gleichgewicht nur ein instabiles sein, und der Apparat, einmal davon abgelenkt, würde nicht dahin zurückkehren, sondern sich immer weiter davon entfernen, und nur in der entgegengesetzten Stellung zur Ruhe kommen, wo der Stab seine natürliche Lage im Raume hat, aber die Aufhängungsfäden einander kreuzen.

Im dritten Falle endlich, wo die beiden Directionskräfte einen Winkel mit einander machen, wird der Conflict dieser beiden Kräfte durch eine Zwischenstellung vermittelt, wobei weder der Stab im Meridian, noch eine gerade Linie durch die untern Anknüpfungspunkte der Fäden der durch die obern parallel ist, und diese Zwischenlage sowohl, als die Kraft, mit welcher der Apparat in derselben zurückgehalten wird, richten sich nach dem statischen Gesetze der Zusammensetzung zweier Kräfte. Man übersieht nun aber zugleich, dass wenn der Apparat Mittel darbietet, die Winkel zwischen den drei in Rede stehenden Stellungen zu messen, das Verhältniss der beiden componirenden Directionskräfte sich berechnen lässt, und dass man folglich auch die magnetische Directionskraft in absolutem Maasse angeben kann, wenn die Directionskraft vermöge der Aufhängungsweise in absolutem Maasse bekannt ist. Unsere Aufgabe ist dann also gelöst. Am vortheilhaftesten ist es übrigens, das Einliegen des Magnetstabes relativ gegen die andern Theile des Apparats so einzurichten, dass jener in der vermittelten Gleichgewichtsstellung nahe einen rechten Winkel mit dem magnetischen Meridian macht, welchem Fall also die Benennung der transversalen Lage vorzugsweise angemessen ist. Theils ist nemlich dann die Ablenkung der Fäden
\SourceLeaf{368}von ihrer Lage in Einer Ebene am grössten, und damit die Berechnung des Resultats am schärfsten, theils hat dann auch die kleine Veränderung der magnetischen Declination vermöge der stündlichen oder zufälligen Variationen auf die Stellung keinen merklichen Einfluss. Dagegen aber afficirt eine jede Veränderung in der Stärke des Erdmagnetismus die Stellung unmittelbar, und lässt sich mit derselben Leichtigkeit, Schnelligkeit und Schärfe sogleich erkennen und messen, wie das Spiel der Veränderungen der Declination am gewöhnlichen Magnetometer.

Die praktische Anwendbarkeit dieser Idee hatte ich schon vor mehreren Jahren durch vorläufige Versuche an einer freilich nur ganz rohen Vorrichtung bestätigt gefunden, wovon auch eine Andeutung in meinem Aufsatze über Erdmagnetismus und Magnetometer [S. 327 d. B.] gegeben ist. Seit kurzem habe ich aber einen vollkommneren Apparat ausführen lassen, und in der Sternwarte an dem Platze, wo sich bisher das Magnetometer mit fünfundzwanzigpfündigem Stabe befand, aufgehängt. Nach den bereits gegebenen Entwickelungen wird sich dieser Apparat kurz beschreiben lassen.

Er ist aufgehängt an zwei 17 Fuss langen Stahldrähten, oder genau zu reden, an einem einzigen, dessen Enden unten an dem Apparat geknüpft sind, während seine Mitte oben über zwei Cylinder geht, die ihn in schicklicher Entfernung (etwa $1\frac{1}{2}$ Zoll) auseinander halten: diese Einrichtung hat zugleich den Vortheil, dass die beiden Stränge von selbst gleiche Spannung haben. Die Aufhängung befindet sich oberhalb der Decke des Saals, und die Drähte hängen frei durch eine kreisrunde $3\frac{1}{2}$ Zoll weite Oeffnung in der Decke. Die Entfernung der Drähte von einander kann sowohl oben als unten nach Bedürfniss weiter oder enger gestellt werden. Der an den Drähten hängende Apparat selbst besteht aus vier Haupttheilen. Der erste, an welchem die Drähte fest sind, ist eine horizontale in Viertelsgrade auf Silber eingetheilte Kreisscheibe, von vier Zoll Durchmesser. Der zweite Theil besteht aus einer auf dem Limbus des Kreises, concentrisch mit diesem drehbaren Alhidade mit zwei Verniers, die einzelne Minuten geben; einer damit fest verbundenen ziemlich starken gegen die Kreisebene senkrechten Stange, und einem daran befindlichen sehr vollkommen kreisrunden Spiegel von $1\frac{1}{2}$ Zoll Durchmesser, in welchem man durch ein 16 Fuss entferntes Fernrohr das Bild eines Stücks einer in einzelne Millimeter eingetheilten unterhalb des Fernrohrs befestigten horizontalen Scale sieht. Auf diese Weise ist
\SourceLeaf{369}also jede Veränderung in der Lage des Kreises zu erkennen und zu messen; kleine Veränderungen unmittelbar mit äusserster Schärfe durch die im Fernrohr sich zeigenden Scalentheile, grössere, indem man damit eine Alhidadenbewegung verbindet und die Verniers abliest. Der dritte Theil ist das unter dem Kreise befindliche Schiffchen, ein doppelter Rahmen, durch welchen der vierte Bestandtheil, ein fünfundzwanzigpfündiger starker Magnetstab gesteckt wird. Dieses Schiffchen ist gleichfalls um das Centrum des Kreises drehbar, und mit zwei auf dem Kreislimbus aufliegenden Verniers versehen, wodurch man die Grösse der Drehung auf die Minute messen kann.

Stellt man nun zuvörderst das Schiffchen so, dass der Apparat einerlei Gleichgewichtslage behauptet, es möge der Magnetstab im Schiffchen liegen, oder ein nicht magnetischer Körper von gleichem Gewicht, so ist dies die erste oder die zweite der vorhin unterschiedenen Hauptlagen, jenachdem der Magnetstab sich dabei in seiner natürlichen, oder in der verkehrten Lage befindet. Die erstere bietet keine besonders wichtige praktische Anwendung dar, und die Brauchbarkeit der zweiten ist an die Bedingung geknüpft, dass die magnetische Directionskraft etwas kleiner sein soll, als die Directionskraft vermöge der Aufhängungsart. Bei dem hiesigen Apparat ist jetzt das Verhältniss dieser Directionskräfte nahe wie 10 zu 11; die resultirende Directionskraft ist also nur der zehnte Theil der magnetischen Directionskraft. Wir haben also hier ein Analogon einer astatischen Magnetnadel, und jede fremde Kraft, die die Richtung einer einfachen Nadel stört, äussert hier eine zehnmal grössere Wirkung, als bei einer Aufhängung an Einem Faden Statt haben würde, und zwar, wie man leicht einsieht, in entgegengesetztem Sinn. Es ist hiedurch also unter anderen die Auflösung einer Aufgabe gegeben, mit welcher man sich früher ohne Erfolg wiederholt beschäftigt hat, nemlich die täglichen und stündlichen Variationen der magnetischen Declination vergrössert darzustellen. Öftere gleichzeitige Beobachtungen dieser Art, an diesem Apparat und am Magnetometer des magnetischen Observatorium haben zwar immer die befriedigendsten Resultate gegeben: inzwischen verliert doch diese Anwendung jetzt von ihrer Wichtigkeit, weil die gewöhnlichen Magnetometer schon die kleinsten Veränderungen mit aller zu wünschenden Schärfe geben, mithin das Bedürfniss einer Vergrösserung jetzt nicht mehr Statt findet.

Diese und andere Anwendungen beim verkehrten Einliegen des Stabes, auf
\SourceLeaf{370}welche ich nachher noch zurückkommen werde, sind jedoch nur als untergeordnete zu betrachten: bei weitem wichtiger ist der Gebrauch des Apparats bei der dritten oder \textit{transversalen} Lage für die Intensität. Wenn man, von der natürlichen Lage ausgehend, durch eine Drehung des Schiffchens den Magnetstab aus dem magnetischen Meridian bringt, so muss sich der ganze Apparat, um zum Gleichgewicht zu kommen, um einen gewissen dem Verhältniss der beiden Directionskräfte entsprechenden Winkel zurückdrehen; die Differenz dieser beiden Winkel wird die Abweichung des Magnetstabes vom magnetischen Meridian in der Gleichgewichtsstellung sein, und man kann es leicht so einrichten, dass diese Abweichung nahe 90 Grad beträgt, wodurch die vorhin bereits angeführten Vortheile erreicht werden. Ganz vorzüglich eignet sich dann aber der Apparat zur Beobachtung der \textit{Änderungen} der Intensität, die sich unmittelbar durch den veränderten Stand kund geben. Dass dabei in Beziehung auf solche Änderungen, die erst nach längerer Zeit erfolgen, mehrere Umstände nicht unberücksichtigt bleiben dürfen, liegt unvermeidlich in der Natur der Sache selbst: namentlich erfordern jene, dass von Zeit zu Zeit durch (bekannte) geeignete Mittel untersucht werde, ob und in welchem Maasse die Stärke des Magnetismus im Stabe sich verändert habe; auch die Temperaturveränderungen kommen in Betracht, einmal insofern sie diese Stärke, und dann auch, insofern sie die Distanz und Länge der Aufhängungsdrähte, und damit die der Aufhängungsart zukommende Directionskraft afficiren. Aber in Beziehung auf die unregelmässigen in kurzen Zeitfristen wechselnden Veränderungen der Intensität leistet nun der Apparat ganz dasselbe, wie das Magnetometer in Beziehung auf ähnliche Änderungen der Declination; auch ist die Beobachtungsart an beiden Apparaten ganz gleich. Die Veränderungen der Intensität erhält man zunächst in Scalentheilen ausgedrückt, die man jedoch leicht auf Bruchtheile der Intensität selbst zurückführen kann. Unter den gegenwärtigen Verhältnissen des Apparats entspricht einem Scalentheile der 22000\textsuperscript{ste} Theil der ganzen Intensität.

Die freilich nur erst eine kurze Zeit umfassenden und nicht sehr zahlreichen bisherigen Erfahrungen an dem Apparat lassen doch schon einige nicht unwichtige Resultate erkennen.

Erstlich deuten die bisherigen Beobachtungen auf regelmässige von der Tageszeit abhängige Änderungen hin, die sich freilich mit unregelmässigen eben so häufig vermengen mögen, wie bei der Declination, und deren sichere Scheidung
\SourceLeaf{371}Jahrelang fortgesetzte Beobachtungen erfordern wird. Wenn ich, nach so wenigen Erfahrungen, wie bisher vorliegen, mehr eine Vermuthung als ein Resultat aussprechen darf, so scheint der regelmässige Gang darin zu bestehen, dass die Intensität in den Vormittagsstunden abnimmt, so jedoch, dass sie schon eine oder zwei Stunden vor dem Mittage ihr Minimum erreicht, und von da an wieder zunimmt. Um doch vorläufig für das quantitative Verhältniss einen Anhaltspunkt zu bekommen, habe ich im August 1837 an 30 Tagen die Stellung Morgens um 10 Uhr und Nachmittags um 3 Uhr aufgezeichnet: das Resultat war, dass an 26 Tagen die Intensität Nachmittags grösser war, und nur an 4 Tagen kleiner, als Vormittags; der mittlere Unterschied betrug 39 Scalentheile, oder etwas mehr als den 600\textsuperscript{sten} Theil der ganzen Intensität. An den meisten jener Tage wurde der Stand des Apparats auch Vormittags um 9 Uhr aufgezeichnet; unter 28 Tagen waren 23, wo die Intensität um diese Stunde noch grösser war, als eine Stunde später, und nur an 5 Tagen fand das Umgekehrte Statt: der mittlere Unterschied betrug hier aber nur $11\frac{1}{2}$ Scalentheile, oder etwas mehr als den 2000\textsuperscript{sten} Theil der ganzen Intensität.

\expandafter\def\csname HistoricalNoteText371\endcsname{Auf ähnliche Art und mit gleichem Erfolge ist später auch in dem Termine vom 13.---14. November an beiden Apparaten beobachtet.}
Zweitens bestätigen mehrere sehr durchgreifende Beobachtungsreihen, dass unregelmässige, zuweilen sehr beträchtliche und in kurzen Zeitintervallen wechselnde Störungen bei der Intensität nicht weniger häufig vorgehen, wie bei der Declination, woran freilich auch an sich nach der Analogie nicht gezweifelt werden konnte. Dreimal schon sind eine beträchtliche Zeit hindurch an diesem Intensitätsapparat und gleichzeitig am Magnetometer des magnetischen Observatorium ununterbrochen fortgesetzte Aufzeichnungen gemacht; am 15. Julius von Morgens 6 Uhr bis Nachmittags 6 Uhr; dann in dem ordentlichen magnetischen Termin vom 29.---30. Julius, endlich in dem ausserordentlichen Termin vom 31. August bis zum 1. September, beidemal 24 Stunden; die Aufzeichnungen geschahen immer von 5 zu 5 Minuten. Graphische Darstellungen der beiden Termine, wo die Curven für die Änderungen sowohl der Intensität als der Declination gezeichnet sind, setzen dies in ein helles Licht. Die beiderseitigen Bewegungen haben zwar, wie sich von selbst versteht, nicht die geringste Ähnlichkeit mit einander; aber sehr bemerklich ist doch, dass wo die Declination stark gestört wurde, meistens auch in der Intensität starke Störungen eintraten\HistoricalNote{371}.
